%%%%%%%%%%%%%%%%%%%%%%%%%%%%%%%%%%%%%%%%%%%%%%%%%
%%%%%%%%%%%% cap: intro %%%%%%%%%%%%%%%%%
%%%%%%%%%%%%%%%%%%%%%%%%%%%%%%%%%%%%%%%%%%%%%%%%%

\chapter{Introduction}\label{cap:intro}

\section{Motivation}\label{sect:motivation}

Large language models are no longer used only as standalone text generators that answer from
their parametric memory. In modern assistants and AI search systems, the model is commonly
connected to external sources such as retrieved web pages, internal knowledge bases, product
catalogs, documentation fragments, and manually assembled context windows. Once that
happens, the final response is shaped not only by the model weights, but also by the content
that is retrieved, selected, ranked, summarized, and injected into the prompt at inference time.

This thesis refers to that steering effect as ``puppeteering'' in a technical sense: an external
actor can influence an LLM response by changing the surrounding data without retraining the
model itself. The term is intentionally provocative, but the phenomenon is not necessarily
malicious. A documentation team may rewrite a page so that a model cites it more accurately.
An internal knowledge engineering team may add supporting context to reduce hallucinations.
At the same time, a vendor may insert strategically crafted text to push a product higher in
LLM-generated recommendations. The same mechanism therefore enables both legitimate
optimization and problematic manipulation.

This problem becomes especially important in generative search. Traditional search engines
mainly return ranked lists of links. Generative engines instead synthesize a direct answer,
usually grounded in a small set of retrieved sources that are quoted, paraphrased, or
prioritized by the model. As a result, content visibility is no longer measured only by whether a
page appears in search results, but also by whether it is surfaced inside the generated answer,
how prominently it appears, and whether its important facts survive summarization. For
content owners, the optimization surface becomes harder to observe. For users, the difference
between well-grounded evidence and carefully engineered phrasing becomes less obvious.

From an engineering perspective, the problem is practical as well as conceptual. Teams need
tools that can measure how pages are represented in generated answers, compare versions of
content, protect factual constraints, and keep a human reviewer in the loop. The repository
developed alongside this thesis addresses that need through \textit{NewGEO}, a prototype
platform for monitoring visibility, generating context-aware recommendations, and validating
content changes before export.

\section{Problem Statement and Research Questions}

The central problem addressed in this thesis is how externally supplied data influences the
content, position, and credibility of LLM responses, and how this influence can be measured
and operationalized in a safe software workflow.

The thesis is guided by the following research questions:
\begin{itemize}
    \item How does retrieved or injected external content affect what an LLM mentions, cites, or recommends?
    \item How can this influence be measured in a repeatable way when direct access to proprietary generative engines is limited?
    \item How can optimization suggestions be generated without losing locked facts or introducing semantic drift?
    \item What system architecture is needed to combine benchmarking, guarded rewriting, human approval, and export in a single workflow?
\end{itemize}

These questions position the thesis at the intersection of retrieval-augmented generation,
generative engine optimization, and trustworthy content engineering. The goal is not to
reverse-engineer a specific commercial model. Instead, the goal is to build an inspectable
platform that makes the influence surface visible, measurable, and reviewable.

\section{Objectives and Proposed Approach}

To answer these questions, the thesis develops a practical software system called
\textit{NewGEO}. The system models the content optimization workflow around a small set of
core entities: projects, tracked pages, content snapshots, context packs, query clusters,
benchmark runs, score snapshots, and recommendations. This makes it possible to move from
an abstract discussion about external influence to an executable end-to-end pipeline.

The proposed approach has five main stages:
\begin{enumerate}
    \item ingest and normalize target pages, either directly or through a crawl-oriented import flow;
    \item define external context through supporting documents, locked facts, required terms, and forbidden claims;
    \item group representative user queries into query clusters and benchmark candidate pages against those queries;
    \item generate context-aware recommendation drafts and estimate their projected visibility lift;
    \item require explicit human approval before exporting optimized content.
\end{enumerate}

This design reflects an important thesis position. External influence on LLM outputs should not
be treated only as an offensive optimization technique. It should also be treated as an
observable and testable software process. Consequently, NewGEO separates measurement,
rewriting, and publication instead of collapsing them into a single opaque step.

\section{Author Contributions}

The main contributions of the thesis are the following:
\begin{itemize}
    \item the design of a modular architecture that separates domain logic, HTTP transport, background processing, and dashboard presentation;
    \item the implementation of an end-to-end backend workflow for page ingestion, context-pack creation, query clustering, benchmark execution, recommendation generation, approval, and export;
    \item a transparent heuristic benchmarking engine that approximates visibility through metrics such as mention rate, citation share, answer position, token share, quote coverage, and run variance;
    \item a constraint-aware recommendation pipeline that preserves locked facts, highlights required terms, ranks supporting documents, and checks for possible semantic drift;
    \item a web dashboard that exposes the workflow as an operational product rather than as a standalone research script;
    \item automated tests for the main service flow, including recommendation generation, benchmarking, crawling, approval, export, and storage backend parity.
\end{itemize}

The current implementation uses deterministic heuristics instead of live commercial LLM APIs.
This is an intentional choice. It keeps the platform reproducible in a local environment and
allows the influence of external data to be studied without hiding core behavior inside opaque
third-party services. At the same time, the architecture leaves clear extension points for future
integration with real connectors and richer retrieval backends.

\section{Thesis Structure}

The remainder of the thesis is organized as follows. Chapter~\ref{chapter:related_work}
reviews the literature on retrieval-augmented generation, generative engine optimization,
strategic manipulation of LLM-based recommendations, and cooperative content rewriting.
The following chapters describe the problem framing for NewGEO, the application
architecture, the implementation details of the platform, and the experimental workflow used
to benchmark and validate the prototype. The thesis concludes with the main findings,
limitations, and directions for future work.
